\input{../tex-inputs/latex-leseansicht-vorspann}

\section[Arthur Schnitzler an Gustav Schwarzkopf, 29. 3. 1900]{L04396 Arthur Schnitzler an Gustav Schwarzkopf, 29. 3. 1900}
\nopagebreak\mylabel{L04396v}
\rehead{ }\normalsize\beginnumbering\briefempfaengerindex{Schwarzkopf, Gustav@\textsc{Schwarzkopf, Gustav}!zzzSchnitzler, Arthur@\emph{von Arthur Schnitzler}!1900-03-291@{29. 3. 1900}|(be}
\toendnotes[C]{\smallbreak\pagebreak[2]}
\correspDesc{Versand  durch Arthur Schnitzler am 29. 3. 1900 in Triest
\newline{}Übermittlung  am 29. 3. 1900 in Triest
\newline{}Zustellung  am 30. 3. 1900 \textbf{Ort fehlend} 
\newline{}Erhalt  durch Gustav Schwarzkopf im Zeitraum [30. 3. 1900 – 3. 4. 1900?] in Wien}\toendnotes[C]{\smallbreak}
\Standort{DLA, A:Schnitzler, HS.1985.1.1897.}
\physDesc{Bildpostkarte, 130 Zeichen
\newline{}Handschrift: Bleistift, deutsche Kurrent
\newline{}Versand: 1) Stempel: »\nobreak{}\oindex{Triest@\textbf{Triest}|pwk}Triest Trieste, 29. \textcolor{gray}{3}. 00, 12M\nobreak{}«.   2) Stempel: »\nobreak{}\oindex{I., Innere Stadt@\textbf{I., Innere Stadt}|pwk}Wien 1/1 1, 30. 3. 00, 11½V–1N, \textcolor{gray}{Bestellt}\nobreak{}«. }\pstart{}{\pb}\textsc{Vienna}\oindex{Wien@\textbf{Wien}|pw}\pend{}\pstart{}\textsc{Herrn Gustav Schwarzkopf}\pend{}\pstart{}Wien\oindex{Wien@\textbf{Wien}|pw}\pend{}\pstart{}\textsc{Tiefer Graben 23}\oindex{Wien@\textbf{Wien}!I., Innere Stadt@\textbf{I., Innere Stadt}!Tiefer Graben 23@\textbf{Tiefer Graben 23}|pw}.\pend{}{\bigskip}
\pstart
           \centering{}{\pb}\textcolor{gray}{\textbf{Saluto da Trieste\oindex{Triest@\textbf{Triest}|pw}.}}\pend
           
\pstart
           \centering{}\textcolor{gray}{\textbf{Canal grande\oindex{Canal Grande di Trieste@\textbf{Canal Grande di Trieste}|pw}.}}\pend
           \vspace{1em}
\pstart
           \noindent{}{\pb}Unerhörtes Wetter! Und ich möchte in einer Stunde nach Raguſa\oindex{Dubrovnik@\textbf{Dubrovnik}|pw} fahren!\pend
           \pstart Herzl Gruß \spacefill\mbox{A.}\pend{}\selectlanguage{ngerman}\endnumbering\briefempfaengerindex{Schwarzkopf, Gustav@\textsc{Schwarzkopf, Gustav}!zzzSchnitzler, Arthur@\emph{von Arthur Schnitzler}!1900-03-291@{29. 3. 1900}|)be}\mylabel{L04396h}
\begin{anhang}
\end{anhang}\newcommand{\dateiname}{L04396}\newcommand{\titel}{Arthur Schnitzler an Gustav Schwarzkopf, 29. 3. 1900}\newcommand{\editorInnen}{Herausgegeben von Jahnke, SelmaMüller, Martin Anton}\input{../tex-inputs/latex-leseansicht-abspann}
